% ---------------------------------------------------------------------
%  Chapitre 7 — Conclusion : bilan, limites globales, perspectives
% ---------------------------------------------------------------------
\chapter{Conclusion}
\label{chap:conclusion}

\section{Bilan}
% [À DÉVELOPPER] Reprendre la question du chap. 1 et y répondre bloc par
% bloc. Résultat marquant du Bloc 1 : seul le RSF (variantes RSI) bat le
% hasard, et cet avantage dépend du régime.
Sur les signaux techniques (Bloc~1), seules deux variantes RSI démontrent
un avantage réel, dépendant du régime de marché. Sur les relations
inter-actifs (Bloc~3), le marché apparaît fortement couplé (facteur
marché à 73~\%) et peu prévisible par son seul passé.

\section{Limites globales}
\begin{limites}
\begin{itemize}
  \item Résultats établis \emph{in-sample}~; une validation hors
        échantillon (\emph{walk-forward}) reste nécessaire avant tout
        usage opérationnel.
  \item Frais et conditions d'exécution modélisés de façon simplifiée.
  \item Puissance statistique élevée sur grands échantillons~: distinguer
        soigneusement significativité et taille d'effet.
\end{itemize}
\end{limites}

\section{Perspectives}
% [À DÉVELOPPER] Finalisation des Blocs 2 et final, validation temporelle,
% mesures de risque (ratios risk-adjusted).
Les perspectives portent sur la finalisation des Blocs~2 et final, la
validation temporelle des signaux et l'ajout de mesures de performance
ajustées du risque.
